\documentclass{resume}

\usepackage[utf8]{inputenc}
\usepackage[T1]{fontenc}
\usepackage[left=0.35in,top=0.20in,right=0.35in,bottom=0.15in]{geometry}
\usepackage{setspace}
\setstretch{1.05}
\emergencystretch=2em
\usepackage{enumitem}
\newcommand{\tab}[1]{\hspace{.2667\textwidth}\rlap{#1}}
\newcommand{\itab}[1]{\hspace{0em}\rlap{#1}}

% Redefine section and name spacing for single-page layout
\renewcommand{\sectionskip}{\vspace{5pt}}
\renewcommand{\sectionlineskip}{\vspace{3pt}}
\renewcommand{\nameskip}{\vspace{3pt}}
\renewcommand{\addressskip}{\vspace{3pt}}

\renewenvironment{rSection}[1]{
  \sectionskip
  {\color{headerblue}\textbf{\MakeUppercase{#1}}}
  \sectionlineskip
  {\color{rulegray}\hrule height 0.6pt}
  \vspace{0pt}
  \begin{list}{}{
    \setlength{\leftmargin}{0em}
    \setlength{\topsep}{0pt}
    \setlength{\partopsep}{0pt}
    \setlength{\parsep}{0pt}
    \setlength{\itemsep}{0pt}
  }
  \item[]
}{
  \end{list}
}

\name{Pham Hung Quoc Viet}
\address{+84 937 404 708 \\ Thu Duc, Ho Chi Minh City \\ \href{mailto:phamviet1415@gmail.com}{phamviet1415@gmail.com} \\ \href{https://github.com/vietpham10205}{github.com/vietpham10205}}

\begin{document}

\begin{rSection}{Objective}
\vspace{-2pt}
{3rd-year student at University of Information Technology -- VNU-HCM (UIT) with a solid foundation in Probability \& Statistics, Time-Series Analysis, and Machine Learning. Experienced in building high-throughput stream processing pipelines, engineering rolling-window statistical features, and benchmarking predictive models using Python and PySpark. Seeking a \textbf{Quantitative Trading / Fintech Intern} position to apply mathematical, statistical, and algorithmic modeling to develop, backtest, and optimize data-driven trading strategies.}
\end{rSection}

\begin{rSection}{Education}
\vspace{-2pt}
University of Information Technology -- VNU-HCM (UIT) \hfill 2023 -- 2027 (expected)\\
\textbf{Bachelor of Information Systems} \hfill GPA: 8.06/10 \\
\textit{Relevant Coursework: Probability \& Statistics, Applied Statistics, Machine Learning, Data Mining, Database Management Systems}
\end{rSection}

\begin{rSection}{Skills}
\vspace{-2pt}
    \renewcommand{\arraystretch}{0.95}
    \begin{tabular}{ @{} >{\bfseries}l @{\hspace{2.5ex}} >{\raggedright\arraybackslash}p{5.4in} @{} }
        Programming           & Python, C++, SQL, Java, PHP \\
        Quantitative          & Probability \& Statistics, Time-Series Analysis, Statistical Modeling \\
        Machine Learning      & Scikit-learn, XGBoost, LightGBM, TensorFlow, Optuna \\
        Data \& Streaming     & Pandas, NumPy, PySpark, Kafka, PostgreSQL \\
        Tools \& Languages    & Git/GitHub, Docker, Jupyter Notebook, IELTS 6.5 (2022) \\
    \end{tabular}
    \renewcommand{\arraystretch}{1.0}
\end{rSection}

\begin{rSection}{Projects}
    \item \textbf{NYC Taxi — Real-time Fleet Intelligence \& Anomaly Detection} \hfill 2026\\
    {\small \textit{PySpark, Kafka, Docker, PostgreSQL, Streamlit, Pandas, Scikit-learn}} \hfill | \href{https://github.com/vietpham10205/real-time-Fleet-Intelligence-anomaly-detection-dashboard}{GitHub}
    {\small
    \begin{itemize}[itemsep=4pt, parsep=0pt, leftmargin=1.2em]
        \item \textbf{High-Frequency Data Ingestion:} Developed a low-latency Kafka pipeline to ingest high-frequency event streams at 200 msg/s, simulating real-time market data architectures.
        \item \textbf{Time-Series Feature Engineering:} Leveraged PySpark Structured Streaming to compute 14 rolling-window statistical features (e.g., moving averages, volatility metrics) on live data streams.
        \item \textbf{Anomaly \& Regime Shift Detection:} Designed an Isolation Forest and Minimum Covariance Determinant (MinCovDet) ensemble with strict voting to identify statistical anomalies and regime shifts.
        \item \textbf{Automated Model Retraining:} Managed continuous data ingestion to PostgreSQL and implemented automated model evaluation and parameter retraining triggered every 25 streaming batches.
    \end{itemize}
    }
    \vspace{15pt}

    \item \textbf{Bank Customer Churn Prediction \& Risk Attribution} \hfill 2026\\
    {\small \textit{Python, Scikit-learn, TensorFlow, XGBoost, LightGBM, SMOTE, SHAP, Optuna, Pandas}} \hfill | \href{https://github.com/ThanhThang5722/Bank-Customer-Churn-Prediction}{GitHub}
    {\small
    \begin{itemize}[itemsep=4pt, parsep=0pt, leftmargin=1.2em]
        \item \textbf{Statistical Preprocessing \& EDA:} Conducted extensive exploratory data analysis; applied Winsorization to mitigate outliers, improving the signal-to-noise ratio in predictive datasets.
        \item \textbf{Class Balancing \& Data Augmentation:} Implemented SMOTE to resolve class imbalance (20.4\% positive rate), statistically augmenting data from 6,400 to 10,192 samples to enhance model sensitivity.
        \item \textbf{Algorithmic Benchmarking:} Trained, optimized, and evaluated 7+ predictive models (XGBoost, LightGBM, Random Forest) utilizing Optuna and GridSearchCV for rigorous hyperparameter tuning.
        \item \textbf{Risk Attribution \& Explainability:} Utilized SHAP and LIME to dissect feature contributions; employed Jaccard similarity analysis to quantify the stability of key risk variables across subsets.
    \end{itemize}
    }
\end{rSection}

\end{document}

